\documentclass[12pt]{article}
\usepackage{amsmath,amssymb,amsfonts}
\usepackage[margin=1in]{geometry}
\usepackage{bm}

\newcommand{\vect}[1]{\mathbf{#1}}

\begin{document}

\section*{Problem 15}

\textbf{Problem:} Prove that the dipole moment of a charge distribution is independent of the origin of coordinates if the total charge is zero.

\bigskip
\textbf{Solution:}

The electric dipole moment of a charge distribution $\rho(\vect{r})$ about an origin $O$ is defined as:
\[
\vect{p} = \int \vect{r}\,\rho(\vect{r})\,d\tau.
\]

Now consider shifting the origin by a vector $\vect{a}$. In the new coordinate system with origin $O'$, the position vector is:
\[
\vect{r}' = \vect{r} - \vect{a}.
\]

The dipole moment about the new origin is:
\begin{align}
\vect{p}' &= \int \vect{r}'\,\rho(\vect{r}')\,d\tau' = \int (\vect{r} - \vect{a})\,\rho(\vect{r})\,d\tau \\
&= \int \vect{r}\,\rho(\vect{r})\,d\tau - \vect{a}\int \rho(\vect{r})\,d\tau \\
&= \vect{p} - \vect{a}\,Q,
\end{align}
where $Q = \int\rho(\vect{r})\,d\tau$ is the total charge.

If the total charge is zero, $Q = 0$, then:
\[
\boxed{\vect{p}' = \vect{p}}.
\]

The dipole moment is independent of the choice of origin when $Q = 0$. \hfill $\blacksquare$

\bigskip
\textit{Remark:} This result is important in physics. For example, a neutral atom or molecule has $Q = 0$, so its dipole moment is an intrinsic property independent of where we place the origin. For a system with $Q \neq 0$, the dipole moment depends on the origin---in the multipole expansion (Example 1.8), the monopole term dominates at large distances, and the dipole term is origin-dependent but physically less significant.

For a discrete charge distribution $\{q_i\}$ at positions $\{\vect{r}_i\}$, the same result holds:
\[
\vect{p}' = \sum_i q_i \vect{r}_i' = \sum_i q_i(\vect{r}_i - \vect{a}) = \vect{p} - \vect{a}\sum_i q_i = \vect{p} \quad \text{if } \sum_i q_i = 0.
\]

\end{document}
